%%%%%%%%%%%%%%%%%%%%%%%%%%%%%%%%%%%%%%%%%%%%%%%%%%%%%%%%%%%%%%%%%%%%%%%%%%%%%%%%%%
% Slide 2 -- section divider
\begin{frame}[fragile]{}
\begin{center}
{\Large Introduction}
\end{center}
\end{frame}

%%%%%%%%%%%%%%%%%%%%%%%%%%%%%%%%%%%%%%%%%%%%%%%%%%%%%%%%%%%%%%%%%%%%%%%%%%%%%%%%%%
% Slide 3 -- Introduction
\begin{frame}[fragile]\frametitle{Introduction}
\begin{itemize}
\item How to model real world objects? --- Design
\item How to put forth ideas in visual manner --- Communication
\item How to verify that design serves the purpose --- Analysis
\item How to get it made? --- Manufacturing
\end{itemize}

All of the above can happen without Computers, but better with Computer assistance.

That's why: Computer Aided \ldots\ (CAx)
\end{frame}

%%%%%%%%%%%%%%%%%%%%%%%%%%%%%%%%%%%%%%%%%%%%%%%%%%%%%%%%%%%%%%%%%%%%%%%%%%%%%%%%%%
% Slide 4 -- History section divider + content
\begin{frame}[fragile]{}
\begin{center}
{\Large History}
\end{center}
\end{frame}

%%%%%%%%%%%%%%%%%%%%%%%%%%%%%%%%%%%%%%%%%%%%%%%%%%%%%%%%%%%%%%%%%%%%%%%%%%%%%%%%%%
\begin{frame}[fragile]\frametitle{History}
\begin{itemize}
\item The first source of CAD resulted from attempts to automate the drafting process.
\item These developments were pioneered by the General Motors Research Laboratories in the early 1960s.
\item Ivan Sutherland's \textit{Sketchpad} (MIT, 1963) established interactive computer graphics as the foundation of modern CAD.
\item CAD became more widely used after 1970 because of technological advancements.
\item CAD allowed users to design products much quicker without the production of an actual product.
\end{itemize}
\end{frame}

%%%%%%%%%%%%%%%%%%%%%%%%%%%%%%%%%%%%%%%%%%%%%%%%%%%%%%%%%%%%%%%%%%%%%%%%%%%%%%%%%%
% Slide 5 -- Evolution of CAD Technology
\begin{frame}[fragile]{Evolution of CAD Technology}
\begin{center}
\includegraphics[width=0.9\linewidth,keepaspectratio]{geomod1}
\end{center}
\end{frame}

%%%%%%%%%%%%%%%%%%%%%%%%%%%%%%%%%%%%%%%%%%%%%%%%%%%%%%%%%%%%%%%%%%%%%%%%%%%%%%%%%%
% Slide 6 -- Manual drafting caption slide
\begin{frame}[fragile]{}
\begin{center}
\includegraphics[width=0.75\linewidth,keepaspectratio]{geomod2}
\end{center}
\vspace{2mm}
\small Manual drafting. Since 1970s: electronic drafting board.
\end{frame}

%%%%%%%%%%%%%%%%%%%%%%%%%%%%%%%%%%%%%%%%%%%%%%%%%%%%%%%%%%%%%%%%%%%%%%%%%%%%%%%%%%
% Slide 7 -- Manual Drafting
\begin{frame}[fragile]\frametitle{Manual Drafting}
\begin{columns}
  \begin{column}{0.55\linewidth}
  \begin{itemize}
  \item 2D representations used to represent 3D objects
    \begin{itemize}
    \item Multi-view drawings (Multiview, Axonometric, Oblique, Perspective)
    \item Pictorials
    \end{itemize}
  \item Standards and conventions developed so 3D object could be built from drawings
  \item Drawings created manually or using 2D CAD
  \item Difficult to visualise, error-prone, time-consuming
  \end{itemize}
  \end{column}
  \begin{column}{0.45\linewidth}
  \begin{center}
  \includegraphics[width=0.9\linewidth,keepaspectratio]{geomod2}
  \end{center}
  \end{column}
\end{columns}
\end{frame}

%%%%%%%%%%%%%%%%%%%%%%%%%%%%%%%%%%%%%%%%%%%%%%%%%%%%%%%%%%%%%%%%%%%%%%%%%%%%%%%%%%
% Slide 8 -- CAD Types
\begin{frame}[fragile]\frametitle{CAD -- Types}
\begin{columns}
  \begin{column}{0.55\linewidth}
  \begin{itemize}
  \item 2D model: Point, line, circular arc, planar curve
  \item 3D model
    \begin{itemize}
    \item Wire frame
    \item Surface
    \item Solid
    \end{itemize}
  \end{itemize}
  \begin{alertblock}{}
  Advantages and Disadvantages of each?
  \end{alertblock}
  \end{column}
  \begin{column}{0.45\linewidth}
  \begin{center}
  \includegraphics[width=0.45\linewidth,keepaspectratio]{CAD2dprofile}
  \includegraphics[width=0.55\linewidth,keepaspectratio]{CADWireframe}
  \includegraphics[width=0.55\linewidth,keepaspectratio]{CADSurface}
  \includegraphics[width=0.55\linewidth,keepaspectratio]{CADSolid}
  \end{center}
  \end{column}
\end{columns}
\end{frame}

%%%%%%%%%%%%%%%%%%%%%%%%%%%%%%%%%%%%%%%%%%%%%%%%%%%%%%%%%%%%%%%%%%%%%%%%%%%%%%%%%%
% Slide 9 -- 2D CAD
\begin{frame}[fragile]\frametitle{2D CAD}
\begin{columns}
  \begin{column}{0.55\linewidth}
  \begin{itemize}
  \item Simply replaces manual drawing
  \item Provides a set of drawing tools to create 2D elements: Lines, circles, arcs, etc.
  \item More accurate, easier changes to drawings
  \item Still no 3D representation of the object
  \item Example: AutoCAD
  \end{itemize}
  \end{column}
  \begin{column}{0.45\linewidth}
  \begin{center}
  
\includegraphics[width=0.9\linewidth,keepaspectratio]{geomod3}
  \end{center}
  \end{column}
\end{columns}
\end{frame}

%%%%%%%%%%%%%%%%%%%%%%%%%%%%%%%%%%%%%%%%%%%%%%%%%%%%%%%%%%%%%%%%%%%%%%%%%%%%%%%%%%
% Slide 10 -- Wire frame caption slide
\begin{frame}[fragile]{}
\begin{center}
\includegraphics[width=0.75\linewidth,keepaspectratio]{geomod4}
\end{center}
\vspace{2mm}
\small Early 1980s: wire frame geometry
\end{frame}

%%%%%%%%%%%%%%%%%%%%%%%%%%%%%%%%%%%%%%%%%%%%%%%%%%%%%%%%%%%%%%%%%%%%%%%%%%%%%%%%%%
% Slide 11 -- 3D Wire frame Modeling
\begin{frame}[fragile]\frametitle{3D Wire frame Modeling}
\begin{columns}
  \begin{column}{0.55\linewidth}
  \begin{itemize}
  \item Geometric entities are lines and curves in 3D
  \item Volume or surfaces of object not defined
  \item Easy to store and display
  \item Hard to interpret --- ambiguous
  \end{itemize}
  \end{column}
  \begin{column}{0.45\linewidth}
  \begin{center}
  \includegraphics[width=0.9\linewidth,keepaspectratio]{geomod4}
  \end{center}
  \end{column}
\end{columns}
\end{frame}

%%%%%%%%%%%%%%%%%%%%%%%%%%%%%%%%%%%%%%%%%%%%%%%%%%%%%%%%%%%%%%%%%%%%%%%%%%%%%%%%%%
% Slide 12 -- Wireframe ambiguity
\begin{frame}[fragile]\frametitle{What is this?}
\begin{center}
Problems with wire frame models

\includegraphics[width=0.8\linewidth,keepaspectratio]{geomod5}
\end{center}
\end{frame}

%%%%%%%%%%%%%%%%%%%%%%%%%%%%%%%%%%%%%%%%%%%%%%%%%%%%%%%%%%%%%%%%%%%%%%%%%%%%%%%%%%
% Slide 13 -- Surface modeling caption slide
\begin{frame}[fragile]{}
\begin{center}
\includegraphics[width=0.75\linewidth,keepaspectratio]{geomod6}
\end{center}
\vspace{2mm}
\small Late 1980s: Surface Modeling
\end{frame}

%%%%%%%%%%%%%%%%%%%%%%%%%%%%%%%%%%%%%%%%%%%%%%%%%%%%%%%%%%%%%%%%%%%%%%%%%%%%%%%%%%
% Slide 14 -- 3D Surface Modeling
\begin{frame}[fragile]\frametitle{3D Surface Modeling}
\begin{columns}
  \begin{column}{0.55\linewidth}
  \begin{itemize}
  \item Models 2D surfaces in 3D space using NURBS and other spline techniques
  \item All points on surface are defined; useful for machining, visualisation, etc.
  \item Surfaces have no thickness; objects have no volume or solid properties
  \item Surfaces may be open
  \end{itemize}
  \end{column}
  \begin{column}{0.45\linewidth}
  \begin{center}
  \includegraphics[width=0.9\linewidth,keepaspectratio]{geomod7}
  \end{center}
  \end{column}
\end{columns}
\end{frame}

%%%%%%%%%%%%%%%%%%%%%%%%%%%%%%%%%%%%%%%%%%%%%%%%%%%%%%%%%%%%%%%%%%%%%%%%%%%%%%%%%%
% Slide 15 -- Surface Model Alias StudioTools
\begin{frame}[fragile]\frametitle{Surface Model --- Alias StudioTools}
\begin{center}
\includegraphics[width=0.7\linewidth,keepaspectratio]{geomod8}
\end{center}
\end{frame}

%%%%%%%%%%%%%%%%%%%%%%%%%%%%%%%%%%%%%%%%%%%%%%%%%%%%%%%%%%%%%%%%%%%%%%%%%%%%%%%%%%
% Slide 16 -- Surface Model Rhino
\begin{frame}[fragile]\frametitle{Surface Model --- Rhino}
\begin{center}
\includegraphics[width=0.7\linewidth,keepaspectratio]{geomod9}
\end{center}
\end{frame}

%%%%%%%%%%%%%%%%%%%%%%%%%%%%%%%%%%%%%%%%%%%%%%%%%%%%%%%%%%%%%%%%%%%%%%%%%%%%%%%%%%
% Slide 17 -- Why draw 3D Models?
\begin{frame}[fragile]\frametitle{Why Draw 3D Models?}
\begin{itemize}
\item 3D models are easier to interpret.
\item Less expensive than building a physical model.
\item 3D models can be altered easily; create more design concepts.
\item 3D models can be used to perform engineering analysis: finite element analysis (stress, deflection, thermal \ldots) and motion analysis.
\item 3D models can be used directly in manufacturing via Computer Numerical Control (CNC).
\end{itemize}
\end{frame}

%%%%%%%%%%%%%%%%%%%%%%%%%%%%%%%%%%%%%%%%%%%%%%%%%%%%%%%%%%%%%%%%%%%%%%%%%%%%%%%%%%
% Slide 18 -- 2D Applications
\begin{frame}[fragile]\frametitle{2D Applications}
\begin{itemize}
\item Drafting --- sketches, architectures, drawings
\item Art --- sketches, painting
\item Electronic layouts, circuit design
\end{itemize}
\begin{center}
\includegraphics[width=0.7\linewidth,keepaspectratio]{geomod10}
\end{center}
\end{frame}

%%%%%%%%%%%%%%%%%%%%%%%%%%%%%%%%%%%%%%%%%%%%%%%%%%%%%%%%%%%%%%%%%%%%%%%%%%%%%%%%%%
% Slide 19 -- 3D Applications
\begin{frame}[fragile]\frametitle{3D Applications}
\begin{itemize}
\item CAD (Computer Aided Design)
\item CAM (Computer Aided Manufacturing)
\item CAE (Computer Aided Engineering) Finite Element Method
\item CG (Computer Graphics)
\end{itemize}
\begin{center}
\includegraphics[width=0.7\linewidth,keepaspectratio]{geomod11}
\end{center}
\end{frame}

%%%%%%%%%%%%%%%%%%%%%%%%%%%%%%%%%%%%%%%%%%%%%%%%%%%%%%%%%%%%%%%%%%%%%%%%%%%%%%%%%%
% Slide 20 -- Basics of FEA
\begin{frame}[fragile]\frametitle{Basics of Finite Element Analysis (FEA)}
\begin{itemize}
\item A complex problem is divided into smaller and simpler problems that can be solved using existing knowledge of mechanics and mathematical tools.
\item Modern mechanical design involves complicated shapes that as a whole cannot be solved by existing mathematical tools.
\item Engineers need FEA to evaluate their designs.
\end{itemize}
\begin{center}
\includegraphics[width=0.65\linewidth,keepaspectratio]{geomod12}
\end{center}
\end{frame}

%%%%%%%%%%%%%%%%%%%%%%%%%%%%%%%%%%%%%%%%%%%%%%%%%%%%%%%%%%%%%%%%%%%%%%%%%%%%%%%%%%
% TODO: Slides 21+ -- export Intro_3DModeling.ppt as PDF and read remaining pages
% to complete this content file.
